\chapter{Discussion}
\label{ch:discussion}

This chapter interprets previously reported results and introduces no new experimental
decisions. Any strengthening of wording remains bounded by the corresponding C01--C11
status; \texttt{descriptive} and \texttt{not\_tested} are not converted into supported
hypotheses here.

The results are not treated as a single ranking of the investigated regimes but as a
sequence of answers to RQ1--RQ3. The central line of the dissertation is the separation
of levels that are easy to conflate: final model behavior, internal gradient and
representation structure, computational-cost structure, mechanistic attribution of
that cost, and the possibility of using intermediate state to control residual
computation. Chapters~\ref{ch:experiments-en} and \ref{ch:results-en} show that moving
between these levels requires a separate estimand and experimental boundary.

\section{RQ1: Proximity to BP Is Not One-dimensional}
\label{sec:discussion-rq1-en}

C01 and C02 together give a coherent but irreducibly multidimensional picture from
Stage~1/2 and Stage~3A. In Stage~1 on FashionMNIST, final macro-F1 values of the
investigated regimes were close, and the registered paired equivalence analysis for
FixedPred and Strict relative to BP lay inside the preregistered \(\pm0.01\) margin.
Stage~2 preserved the registered final-quality surface while the runtime ratio relative
to BP changed. Quality and the price of obtaining it are therefore already distinct.

Stage~3A explains why practical equivalence of a final metric cannot substitute for
internal equivalence. FixedPred gradient direction remains close to BP in the studied
domain, but early-gradient norm is strongly reduced. Strict shows more visible
disagreement in hidden-gradient direction and representation measures. C02 therefore
concerns concrete coordinates of internal state---gradient direction, gradient scale,
and representation similarity---rather than a general claim of ``algorithm
equivalence''.

This result is consistent with the literature boundary in Chapter~\ref{ch:related-work}.
Theoretical PC/BP relations depend on concrete dynamics, initialization, and update
ordering \cite{whittington2017approximation,song2020exact,millidge2022arbitrary,
rosenbaum2022relationship,rosenbaum2025correction}. The empirical contribution here is
an executable counterpart to that distinction: even when final metrics do not separate
regimes beyond a registered operational boundary, gradient and representation
diagnostics can distinguish their internal trajectories.

The answer to RQ1 is therefore conditional. Some regimes and metrics show proximity to
BP in the studied domain, but the domain of that proximity depends on analytical
level. Exact serves as numerical and technical control. FixedPred has closer observed
agreement with BP than Strict in gradient direction and representations in registered
Stage~3A, but it does not preserve early-gradient scale. These observations are not
collapsed to one scalar score and are not transferred beyond the studied
implementations, datasets, and model seeds.

\section{RQ2: From Runtime Difference to Cost Localization and Attribution}
\label{sec:discussion-rq2-en}

RQ2 follows from a limitation of RQ1: an observed runtime difference does not reveal
which part of computation creates it or whether that part can be changed while
preserving registered behavior.

B0 localizes the principal device-time region to \texttt{state\_inference}. On the
frozen ROCm/float32 matrix, Strict has higher device time and peak allocated memory than
FixedPred, and \texttt{state\_inference} dominates device time for both regimes. The
independent B0 unit is \texttt{model\_seed}, with three seeds per configuration;
runtime and memory ratios are therefore interpreted as descriptive engineering
estimates on the registered profiling surface rather than highly precise population
estimates. C03 is localization of observed cost within a concrete execution environment,
not independent proof of a causal mechanism or a universal cost claim for every PC
implementation. Mechanistic attribution is tested only in the following SI-MA0/SI-MA1
stages.

The SI-MA0 $\rightarrow$ SI-MA1 sequence has methodological importance of its own. The
initial SI-MA0 mechanistic cost-attribution model fails the registered COST-MA0 gate.
That negative result is preserved as C04 and is not replaced by later calibration. C05
(\texttt{supported}) records the independent SI-MA1 calibration, which asks the
narrower question of whether observer-cost calibration passes CAL-COST-MA1. It does in
the registered domain. The signed residual is interpreted as over-closure relative to
the accounting equation, not negative physical duration. The residual is therefore not
converted into artificial ``budget'' for later stages.

B1 and B2 add another separation: functional admissibility is tested before comparative
profiling. Passage of EQ-B1 and EQ-B2 supports C06 only for the corresponding exact
candidates. The 288-cell matched profiling surface then completes with no failed or
repeated runs. The registered engineering-continuation screen applied to four groups of
16 configurations yields 0/16 qualified configurations in every group; both candidates
receive \texttt{reject\_or\_revise}. C07 records this result as \texttt{descriptive}.

Stage~3B therefore exposes a difference among admission levels: numerical and trajectory
equivalence of B1/B2 does not guarantee passage of the registered multi-component
resource-continuation screen. This is not a claim that
\texttt{stage2\_baseline} is superior. The publication boundary explicitly forbids a
superiority claim, and \texttt{reject\_or\_revise} concerns continuation of the studied
candidates on the registered engineering matrix only.

Taken together, RQ2 receives a sequential localization rather than one ``main source of
cost''. B0 identifies \texttt{state\_inference}; cost attribution then keeps the failed
initial model distinct from observer calibration; exact candidates next pass their own
numerical and trajectory gates but neither passes the separate engineering-continuation
screen on the studied matrix. This organization reduces the risk of conflating
correctness, mechanistic attribution, and engineering acceptability: numerical
equivalence does not imply resource admissibility under the registered screen.

Recent literature reinforces that PC efficiency and scaling are independent research
axes. iPC changes temporal update order for stability and efficiency
\cite{salvatori2024stable}; PCX systematizes scalability benchmarking
\cite{pinchetti2025benchmarking}; $\mu$PC addresses training of networks deeper than
one hundred layers \cite{innocenti2025mupc}; and ePC examines digital cost and signal
propagation with depth \cite{goemaere2026epc}. These works do not explain the B0 or
Stage~3A mechanisms in the present study. In particular, the reduced early-gradient
norm of FixedPred is not identified with the signal attenuation discussed in ePC; such
a causal comparison would require a separate experiment sharing a mechanistic
diagnostic.

\section{RQ3: Sufficiency, Recognizability, and Economics}
\label{sec:discussion-rq3-en}

The final QWake-FP question is a bounded empirical test of the selective PC-TREF case,
not proof of full sufficiency on the registered domain
$E_I|_\Omega\subseteq E_R|_\Omega$. It separates three levels: presence of
preterminal records whose states belong to \texttt{EARLY\_ADMISSIBLE}; ability to
select a non-zero subset of that class on the frozen calibration surface with zero
observed dangerous accepts; and economic viability of the decision mechanism. The
levels are logically related but not interchangeable.

Structural reconciliation of the temporal surface identifies exactly which accepted
records support the first two levels. The frozen C1 request defines 108 trajectories,
and the full surface contains seven epochs \(t=0,\ldots,6\) per trajectory. The
highest-coverage zero-danger rule, \texttt{compute\_step >= 5}, accepts two epochs: 108
records at \(t=5\), where one sweep remains, and 108 terminal-boundary records at
\(t=6\). Because the rule has zero registered dangerous accepts, the 108 step-5 records
form an explicitly preterminal accepted subset with zero observed dangerous accepts.
The presence of preterminal states for which the early action is admissible is therefore
supported without relying on the terminal boundary.

Within the frozen family of 2,625 scalar rules, non-zero calibration recognizability
with zero observed dangerous accepts exists: 265 rules have zero dangerous accepts and
264 of those also have non-zero coverage. The highest-coverage rule's registered
coverage remains 216/756 (about 28.57\%), but that quantity refers to the full C1
surface and includes 108 terminal records; it is not the fraction of preterminal early
actions. These findings support C08 in two bounded senses: preterminal records exist on
the studied surface for which the registered analytic completion is required-response
equivalent to the exact suffix, and the frozen family contains a pre-action signal that
selectively identifies a non-zero preterminal subset with zero observed dangerous
accepts. In PC-TREF terms, the observed finite manifestation on the preterminal surface
is
\[
  \{\rho\in\mathcal R_{\mathrm{pre}}:\pi(x(\rho))=a_{\mathrm{early}}\}
  \subseteq
  \{\rho\in\mathcal R_{\mathrm{pre}}:x(\rho)\in S_{a_{\mathrm{early}}}\},
\]
not proof of full sufficiency of \(\phi_I\) on \(\Omega\). It also establishes no
statistical safety beyond the 756 calibration records.

The economic level receives a different answer. None of the 2,625 rules simultaneously
satisfies zero dangerous accepts, non-zero coverage, and positive aggregate net saving
under frozen full decision-cost accounting. C09 is therefore rejected within the exact
C2 estimand. This is a bounded negative result: it excludes such a rule only within the
studied family, interpreter, evidence surface, and accounting model, not across all
possible recognizers or implementations.

C08 and C09 must therefore not be merged. If zero-danger recognizability had been absent
even on the calibration surface, there would be no empirical basis here for a later
engineering question about a cheaper recognizer. If C09 had been supported, C2 could
have frozen a rule and opened confirmatory C3. The observed result lies between these
extremes: a non-zero zero-danger calibration region exists, but the frozen accounting
makes every evaluated rule economically non-beneficial. Bounded informational
viability is supported; economic viability under this accounting is not.

The result must also be mapped correctly to theory. It is not a negative result for
PC-TREF: the non-zero calibration region with zero observed dangerous accepts shows
that, within the studied family and surface, the diagnostic representation contains
information sufficient for a subset of decisions. It does not falsify PC-CATM because
C2 does not test universal correctness of that mechanistic model. Nor does C2 provide
positive evidence that PC-CATM features are superior: the highest-coverage zero-danger
rule uses only \texttt{compute\_step >= 5}, not channel activity, resulting correction,
cancellation coefficient, or transport measures. The result therefore does not show
that PC-CATM features are necessary or better than a simple trajectory-progress
indicator. Moreover, \texttt{compute\_step >= 5} is a fixed-prefix temporal rule and
does not demonstrate input-content-dependent adaptivity. PC-CATM remains a
mechanistically motivated diagnostic feature system; direct confirmatory comparison of
its NCZ/ECZ/TNZ and transport/cancellation features was not performed. Their comparative
value for compute control requires a separately preregistered feature comparison.

\section{Observer-cost Dominance and Its Interpretation Boundary}
\label{sec:discussion-observer-cost-en}

Cost decomposition of the highest-coverage zero-danger rule localizes the barrier in C2.
Under frozen accounting, aggregate full decision cost is about 15,096 seconds, while
the registered \texttt{remaining\_suffix\_ns} sum over accepted non-dangerous records
is about 4.98 seconds. Approximately 99.887\% of full decision cost belongs to the
observer category. Break-even with the same registered suffix sum would require about
99.967\% reduction of the accounted cost.

The 4.98-second quantity is not identified with pure time of avoidable iterative sweeps.
The accepted set includes the terminal boundary \(t=6\), where remaining sweeps are
zero even though measured \texttt{remaining\_suffix\_ns} may include service work of
exact completion and accounting. The structural reconciliation therefore does not
remove terminal records and does not recompute aggregate net saving or the original C2
estimand.

These quantities describe the frozen full decision-cost surface, not measured runtime
of the single threshold comparison \texttt{compute\_step >= 5}. C2 assigns every rule
the same per-record cost sequence derived from measured C1 edges. This provides a
common outcome-independent accounting model: after seeing rule results, a preferred
rule cannot be assigned a cheaper cost model. But the same design bounds the estimand.
C09 is rejected under this registered full accounting, not under an unmeasured marginal
runtime model for a lightweight recognizer.

Consequently, the ratio of full decision cost to the registered suffix-time sum cannot
be interpreted as the runtime ratio of a lightweight recognizer. C10 is therefore
\texttt{not\_tested}. Testing it requires a new experiment in which the marginal
observation/admission path is materialized, independently measured, and preregistered
under the same dangerous-accept criterion. Recomputing old C2 post hoc with a lower cost
would change the estimand after the outcome is known.

\section{Significance of Negative Results}
\label{sec:discussion-negative-results-en}

The research line contains at least two scientifically meaningful negative tests:
SI-MA0 and the QWake-FP C2 economic test. They serve different roles but jointly show
why preservation of negative results matters for mechanistic interpretation.

SI-MA0 bounds the first mechanistic cost-attribution model. Later SI-MA1 calibration
does not turn COST-MA0 into a passed gate; it tests a new narrower construction.
Likewise, the bounded-negative C2 does not become positive merely because cost
decomposition reveals observer dominance. Decomposition explains \emph{where} the
accounting barrier lies but does not change the original eligibility condition.

This logic prevents two forms of post-hoc interpretation. A failed gate cannot be
silently replaced by a new threshold and called the same test, and a diagnostic result
cannot retroactively authorize a stage that the preregistered protocol opened only
after a positive predecessor receipt. Absence of C3 is therefore part of the result:
C11 remains \texttt{not\_tested} because C2 created no rule-freeze receipt.

\section{QWake-PC/QWake-FP Relative to Adaptive Computation and Selective Prediction}
\label{sec:discussion-related-qwake-en}

QWake-PC is close to adaptive-computation and early-exit literature in its economic
formulation: an extra decision mechanism is useful only when its cost is compensated by
avoided computation \cite{graves2016adaptive,bolukbasi2017adaptive}. For PC
specifically, Zniber et al. implement dynamic early exit from PC cycles using separate
classifiers and a confidence threshold \cite{zniber2025earlyexit}. The dangerous-accept
test also places QWake near selective prediction and risk-aware early exit, where
coverage is considered jointly with risk
\cite{geifman2019selectivenet,jazbec2024riskcontrol,valade2025eero}.

Direct identification would nevertheless be incorrect. In EE-PCN and typical
multi-exit networks, an intermediate prediction head is part of the model and may be
trained specifically for early exit. QWake-FP C2 uses a bounded representation of an
already existing PC iterative computation, available before action, while admissibility
of the registered analytic completion is labeled by comparing its required response
with the post-action exact-suffix required response. QWake-FP novelty is therefore not
formulated as the first application of early exit to PC. The C2 dangerous-accept
criterion is also a bounded empirical condition, not a statistical risk-control
risk bound imported from the literature.

The QWake-FP contribution in this dissertation is instead the experimental separation
of three questions on one sealed surface: is the registered analytic replacement
admissible relative to the exact required response; is a pre-action signal available
that selects it with zero observed dangerous accepts; and does the decision pay under
the preregistered accounting model? The answers differ, and that difference defines the
next research direction.

\section{Threats to Validity}
\label{sec:discussion-validity-en}

\subsection{Construct Validity}

``Proximity to BP'', ``cost'', and ``safety'' are represented by concrete operational
definitions. Proximity is evaluated separately through final metrics, gradient
direction and magnitude, and representation similarity; none substitutes for all the
others.

In Stage~3B, cost is represented by registered device time, memory, observation, and
accounting components. In QWake, safety means absence of dangerous accepts on sealed
calibration evidence relative to a concrete post-action oracle. This is narrower than
a statistical population-risk bound.

\subsection{Internal Validity}

Versioned protocols, frozen configurations, exact source/image identities, checksums,
and admission gates reduce the risk of silent condition drift between stages. Matched
designs further reduce mismatched execution conditions. The program nevertheless
evolves over time: Stage~1/2, Stage~3A, Stage~3B, and QWake use different local
estimands and instrumentation surfaces. Cross-stage relations are therefore interpreted
as sequential narrowing, not as one monolithic experiment with an invariant measurement
system.

For QWake-FP, an important protection is separation of pre-action rule input from the
post-action oracle. C2 is nevertheless a bounded calibration search, not independent
confirmation. The absence of C3 limits generalization of zero-danger recognizability
observed on the calibration surface.

\subsection{Statistical Validity}

Primary Stage~1/2 and Stage~3A results use ten independently trained models. B0 has a
narrower replication base: \texttt{model\_seed} is the independent unit with three
values per configuration, so its time and memory ratios are descriptive engineering
estimates. Layer observations within one model are not independent replications. This
prevents pseudoreplication but leaves limited power for small effects and limits
precision across model seeds.

Stage~1 separates difference testing from equivalence testing; failure to detect a
difference is not evidence of equivalence. QWake-FP C2 is not a population-risk
estimate: zero dangerous accepts refer to 756 sealed calibration records and the frozen
rule family.

\subsection{External Validity}

The experimental domain is bounded to Torch2PC, FashionMNIST and MNIST,
\code{lenet\_classic}, the concrete Exact/FixedPred/Strict implementations, the
float32/ROCm profiling environment, and registered execution matrices. Other
architectures, datasets, hardware, numerical precision, or PC dynamics can exhibit
different gradient, representation, and cost behavior.

The QWake-FP economic result has an especially narrow transfer boundary. Observer-cost
dominance belongs to the C1/C2 measurement architecture. Another runtime architecture
may change the ratio between observation cost and avoidable residual computation. C09
must therefore travel only with its applicability boundary.

\subsection{Reproducibility and Artifact Provenance}

The repository stores tracked evidence, frozen specifications, and thesis-facing
summaries with SHA-256 bindings to sources. CI checks selected aggregates before LaTeX
generation. This reduces manual text/data drift but does not imply that every external
reader can rerun every historical scientific execution without the original hardware
environment and private or locally retained forensic artifacts.

Scientific-result reproducibility and byte identity of a PDF are also distinct. The
thesis CI identifies each concrete PDF by SHA-256; that checksum does not by itself
establish byte determinism across independent builds.

\section{What the Current Work Does Not Establish}
\label{sec:discussion-nonclaims-en}

To prevent claim expansion beyond the evidence surface, the dissertation explicitly
does \emph{not} establish:
\begin{itemize}
  \item universal equivalence of PC and BP;
  \item universal superiority of FixedPred, Strict, B1, or B2;
  \item population-level safety from zero dangerous accepts in C2;
  \item impossibility of a recognizer outside the frozen family of 2,625 rules;
  \item marginal runtime cost of a simple threshold recognizer from the observer-dominated full decision cost;
  \item execution or implicit passage of confirmatory C3.
\end{itemize}
These are part of the claim contract, not caveats added after the result. They define
the boundary within which RQ1--RQ3 remain traceable to the registered evidence surface.

\section{Directions for Future Work}
\label{sec:discussion-future-en}

The first follow-up question follows directly from C08, C09, and C10: can a lightweight
recognizer preserve a preregistered dangerous-accept criterion on a new confirmatory
surface at far lower marginal observation/decision cost? This requires a new
preregistered estimand measuring the concrete recognizer's actual extra path. It must
not revalue the old C2; it needs a new identifier, new cost surface, and its own
break-even criterion.

A second direction is independent confirmation of recognizability on a new evidence
surface. It is meaningful only after a new eligible rule or recognizer exists and is
not continuation of the closed original C3 chain.

A third direction is broader external validity: additional datasets, architectures,
depths, numerical-precision regimes, and hardware. The separation among behavioral
proximity, internal mechanism, and cost should be preserved rather than transferring a
claim from one analytical level to another.

Mechanistic work can also profile alternative observer architectures and cheaper
pre-action state representations. The scientifically relevant result would not be lower
overhead alone, but preservation of the preregistered dangerous-accept condition under
independent measurement.

\section{Discussion Summary}

The main interpretation of the completed program is that different notions of
viability receive different answers. Behavioral proximity to BP is observed for some
regimes and metrics but does not mean internal identity. Profiling localizes
substantial cost, while exact alternative organizations, despite passing equivalence
gates, receive \texttt{reject\_or\_revise} on a separate engineering-continuation
screen; this is not baseline superiority. Finally, QWake finds bounded calibration
recognizability with zero observed dangerous accepts, but frozen full decision-cost
accounting admits no rule that combines the same condition with positive economic
benefit.

This is not an unfinished experimental program. The stopping rule worked as designed:
absence of an eligible C2 rule closed the confirmatory chain, preserved C3 as
\texttt{not\_tested}, and localized the next scientific question in a new estimand---
the marginal cost of a recognizer preserving the preregistered dangerous-accept
criterion.
